\subsection{Which Fourier features help on real crystals?}\label{sec:revision-ablation}
Two eight-site cells demonstrate a loss of phase information in our binary fingerprint (Fig.~\ref{fig:revision-collision}). They share their lattice, composition, all 13 sampled species powers and cross-powers, and all six mixed triple products. Their oxygen motifs have different circular gap multisets and cannot be related by translation or reversal. One signed triple along a single axis distinguishes them exactly. Sections~S3 and S9 give the derivation and its dependence on the lattice basis.

\begin{inplace}{figure}
 \centering\includegraphics[width=.92\linewidth]{revision-fourier-collision.pdf}
 \caption{\textbf{A phase distinction missed by the trained fingerprint.}
 The horizontal axis in panels a and b covers one cell period. Oxygen sites lie at $y=0$ and silicon sites at $y=1/3$. Both motifs use $z=0$ throughout. Open grey sites mark the remaining twelfths of the first lattice vector. The underlying cyclic pair appears in Mandereau et al.\ \cite{zrelation}, Fig.~1. Panel c shows the coincident oxygen power spectra for axial harmonics 0 to 6. Panel d shows the opposite signs of an additional, unscaled collinear triple. The original fingerprint agrees to $4.2\times10^{-17}$ componentwise in the two cells.}
 \label{fig:revision-collision}
\end{inplace}

To test whether this distinction helps predict held-out crystals, we ran four binary feature sets under identical training conditions. The 58-component pair descriptor contains the cell, the species powers and the cross-power. The original descriptor adds twelve real components from its six complex triples. A third arm appends the signed collinear channel to the original. A fourth instead duplicates an existing pair feature and rescales it to produce the same mean increment in squared distance on training-only real/generated pairs. All arms use the same three initializations, 1,200 updates, sampled reference indices and independent noise draws. Retraining the original reproduced its archived parameters exactly at all three seeds.

\begin{center}\small
\begin{tabular}{lrrrr}
\hline
 & \multicolumn{2}{c}{Environment energy score $\downarrow$}
 & \multicolumn{2}{c}{Recovery in five draws (\%) $\uparrow$}\\
Training features & Validation & Test & Validation & Test\\\hline
Pairs only (58) & 3.4343 & 3.4282 & 46.61 & 43.23\\
Original (70) & 3.4711 & 3.4693 & 41.93 & 43.75\\
Original + signed triple (71) & 3.4686 & 3.4690 & 42.97 & 42.71\\
Original + reweighted pair (71) & 3.4705 & 3.4706 & 42.97 & 43.49\\\hline
\end{tabular}
\end{center}
The pair-only arm had the lowest validation energy score. On test it improved the score by $0.0410$ relative to the original (paired 95\% case-bootstrap interval $[0.0354,0.0466]$). Its five-draw recovery differed by $-0.52$ percentage points ($[-3.65,2.86]$). The added signed channel repaired the analytic collision, but its test improvements did not reproduce the small validation energy change. Its test score difference from the original was $-0.0002$ ($[-0.0025,0.0020]$) and its recovery difference was $-1.04$ percentage points ($[-3.12,1.04]$). These results distinguish the information a feature can represent from the structural accuracy learned at a fixed training budget.
